\subsection{Background of the expert policies}
\label{sec:expert-background}

\qquad The expert used in this study is adapted from Team 4's controller in
the SuperTuxKart Bots race project. That project asks student teams to build
agents that win races using simulator telemetry. Its agents are hand-written
controllers: the term \emph{expert} refers to specialized driving rules, rather
than a policy learned through reinforcement learning.
Some teams provide Optuna scripts for tuning controller parameters. For example,
Team 2's objective minimizes race steps, while Team 5 minimizes the mean number
of race steps plus 0.2 times its standard deviation. These scripts document
tuning objectives, but do not establish which optimization runs produced the
submitted configurations.\footnote{Sources:
\href{https://github.com/osigaud/pystk_project/blob/ddf2cb6573850f890bba26d4e5f82888d9dac82a/README.md}{project brief},
\href{https://github.com/osigaud/pystk_project/blob/ddf2cb6573850f890bba26d4e5f82888d9dac82a/src/agents/team2/optimise_kart.py}{Team 2 tuning}, and
\href{https://github.com/osigaud/pystk_project/blob/ddf2cb6573850f890bba26d4e5f82888d9dac82a/src/agents/team5/optimization.py}{Team 5 tuning}.}

\subsubsection{Racing objective and items}

\qquad The intended task is racing performance, rather than collecting as many
bonus boxes as possible. The upstream environment defines the per-step reward
as
\[
r_t = \frac{d_t-d_{t-1}}{10}
    + \left(1-\frac{p_t}{N}\right)(3+7f_t)
    - 0.1 + 10f_t,
\]
where \(d_t\) is nonnegative overall distance, \(p_t\) is race position,
\(N\) is the number of karts, and \(f_t\) indicates that the kart has finished.
There is no direct reward for collecting boxes. Items can improve progress,
position, or finishing time indirectly. Moreover, this environment reward
should not be confused with the objective of a hand-written controller:
Team 4 selects actions through rules, without learning from this reward.
Maximizing the discounted return is also not generally equivalent to minimizing
lap time.\footnote{Source:
\href{https://github.com/osigaud/pystk_project/blob/ddf2cb6573850f890bba26d4e5f82888d9dac82a/src/pystk2_gymnasium/envs.py}{upstream reward implementation}.}

The original Team 4 controller does account for items. Its coordinator combines
road following with recovery, avoidance of bananas and bubblegum, and avoidance
of other karts. The driving module also uses collected power-ups and activates
nitro when the steering angle and upcoming curvature allow it. These mechanisms
support racing; its road-following target does not explicitly seek bonus
boxes.\footnote{Sources:
\href{https://github.com/osigaud/pystk_project/blob/ddf2cb6573850f890bba26d4e5f82888d9dac82a/src/agents/team4/agent4.py}{Team 4 coordinator},
\href{https://github.com/osigaud/pystk_project/blob/ddf2cb6573850f890bba26d4e5f82888d9dac82a/src/agents/team4/AgentPilot.py}{driving module}, and
\href{https://github.com/osigaud/pystk_project/blob/ddf2cb6573850f890bba26d4e5f82888d9dac82a/src/agents/team4/AgentItems.py}{power-up rules}.}

\subsubsection{Track independence and the solo adaptation}

\qquad Team 4's driving core uses pure pursuit: it steers toward an upcoming
road point expressed relative to the kart, and reduces throttle as local
curvature increases. It uses a shared configuration rather than a track-name
lookup or a memorized route. It can therefore operate on different tracks
without retraining. This is track independence by construction, not a guarantee
of successful driving on every layout; fixed gains and lookahead can still
fail on difficult turns.\footnote{Sources:
\href{https://github.com/osigaud/pystk_project/blob/ddf2cb6573850f890bba26d4e5f82888d9dac82a/src/agents/team4/steering.py}{pure-pursuit steering} and
\href{https://github.com/osigaud/pystk_project/blob/ddf2cb6573850f890bba26d4e5f82888d9dac82a/src/agents/team4/speed.py}{curvature-based throttle}.}

\qquad Our solo-driving expert retains this road-following core, with lookahead
node 2, wheelbase parameter 2.5, steering gain 7, and curvature-based throttle.
It removes item, nitro, opponent, and explicit recovery logic, so each expert
label is a deterministic function of the current observation. Expert and learned
policies share the same controls and native automatic rescue behavior. The
expert is a demonstrator rather than a claimed minimum-time controller.
Qualification found that the fast variant failed on Zen Garden, which was
therefore excluded from training and transfer evaluation. This limitation
separates a controller's ability to accept different tracks from evidence that
it drives them reliably.

